The proof method of Theorem \ref{1} is similar to the proof of the sharp Lieb--Thirring inequa\-lity~\eqref{LTSch} for a one-dimensional Schr\"odinger operator in the case $\gamma=1/2$ as presented in \cite{HLW}. It~relies on a property of convolutions of the resolvent kernels of the operator under consideration. Such a semigroup property was also recently established for Jacobi operators where it was again used to prove sharp Lieb--Thirring type inequalities \cite{LLS}. With a different proof (not using the convolution property) the sharp inequalities for the Schr\"odinger operator and the Jacobi operator were first obtained in \cite{HLT} and in \cite{HS}, respectively. Despite formal similarity to the case of Jacobi operators, it is still surprising that the proof method works for functional difference operators $W_V(b)$. These operators could be considered as differential operators of infinite order since the symbol $\cosh(2\pi b k)$ can be written as an infinite Taylor series of symbols of even degree w.r.t.\ the variable $k$.

\section{Free resolvent}

Since $W_0(b)\ge 2$ we conclude that $W_0(b)-\lambda$ is an invertible operator for $\lambda<2$. Let
$\lambda=-2\cos(\omega)$ with $\omega\in[0,\pi]$ if $\lambda\in[-2,2]$ and $\omega\in\ii[0,\infty)$ if $\lambda<-2$. Then in Fourier space the inverse of $W_0(b)-\lambda$ is given by the multiplication operator $(2\cosh(2\pi b k)+2\cos(\omega))^{-1}$.

Applying the inverse Fourier transform $\mathcal{F}^{-1}$ to $(2\cosh(2\pi b k)+2\cos(\omega))^{-1}$ we find the kernel of the free resolvent $G_\lambda=(W_0(b)-\lambda)^{-1}$ that is
\begin{gather}\label{rezolvent}
G_\lambda(x,y) = G_\lambda(x-y)=\frac{1}{2b\sin \omega} \frac{\sinh \big(\frac{\omega}{b} (x-y)\big)}{\sinh \big(\frac{\pi}{b} (x-y)\big)}.
\end{gather}

\begin{Remark}
Note that $G_\lambda(x-y)$ is an even and positive kernel for $\omega\in[0,\pi]$ and it becomes oscillating if
$\omega\in \ii(0,\infty)$.
This fact is one of the reasons why we are able to study Lieb--Thirring inequalities only for the eigenvalues $\lambda_j \in [-2,2]$. This interval contains all of the discrete spectrum if the potential $V$ is ``small'' enough. However, if $V$ generates eigenvalues lying in $(-\infty,-2)$, then the oscillating property of the Green's function prevents us from obtaining the desired inequality for all eigenvalues.
\end{Remark}

Note that the value of $G_\lambda$ on the diagonal $x=y$ takes the form
\begin{gather}\label{G0}
G_\lambda(0)=\frac{1}{2\pi b}\frac{\omega}{\sin\omega}
\end{gather}
and we can see the relation between the right-hand side of \eqref{G0} and the expression in the left-hand side of \eqref{main}.
Due to our parameterisation of the spectral parameter, the convergence $\lambda\to 2^-$ implies $\omega\to \pi^-$ and thus
\begin{gather*}
G_\lambda(0) \sim \frac{1}{2b} \frac{1}{\sqrt{1-\cos^2\omega}} \sim \frac{1}{2b} \frac{1}{\sqrt{2-\lambda}}
\qquad {\rm as} \quad \lambda\to 2^-.
\end{gather*}
If $\lambda\to-\infty$, then $\omega \to \ii\infty$ and
\begin{gather*}
G_\lambda(0) \sim \frac{1}{\pi b} |\lambda|^{-1} \log |\lambda|.
\end{gather*}
In \cite{FT2} L.~Faddeev and L.A.~Takhtajan studied the resolvent in a slightly different form
\begin{gather*}
G_\lambda(x,y) = \frac{\sigma}{\sinh \big(\frac{\pi \ii \varkappa}{\sigma}\big) } \bigg( \frac{\e^{-2\pi \ii \varkappa(x-y)}}{1-\e^{- 4 \pi \ii \sigma (x-y)}} +
\frac{\e^{2\pi \ii \varkappa(x-y)}}{1-\e^{4 \pi \ii \sigma (x-y)}} \bigg),
\end{gather*}
which coincides with \eqref{rezolvent} with $\sigma = \ii/2b$, $\lambda = 2\cosh(2b\pi\varkappa)$ and $\varkappa = \frac{\omega - \pi}{2\pi \ii b}$.
It was pointed out that the free resolvent can be written using the analogues of the Jost solutions
\begin{gather*}
f_-(x,\varkappa) = \e^{-2\pi \ii \varkappa x} \qquad {\rm and} \qquad f_+(x,\varkappa) = \e^{2\pi \ii \varkappa x}
\end{gather*}
that appear in the theory of one-dimensional Schr\"odinger operators. Namely
\begin{gather*}
G_\lambda(x-y) = \frac{2\sigma}{C(f_-,f_+)(\varkappa)}
\bigg(\frac{f_-(x, \varkappa) f_+(y, \varkappa)}{1 - \e^{\frac{\pi \ii}{\sigma'} (x-y)}} +
\frac{f_-(y, \varkappa) f_+(x, \varkappa)}{1 - \e^{-\frac{\pi \ii}{\sigma'} (x-y)}} \bigg),
\end{gather*}
where $\sigma' \sigma = -1/4$ and where $C(f,g)$ is the so-called Casorati determinant (a difference analogue of the Wronskian) of the solutions of the functional-difference equation
\begin{gather*}
C(f,g)(x,\varkappa) = f(x+ 2\sigma', \varkappa)g(x,\varkappa) - f(x,\varkappa) g(x + 2\sigma', \varkappa).
\end{gather*}
For the Jost solutions $C(f_-,f_+)(x,\varkappa) = 2\sinh\big(\frac{\pi \ii\kappa}{\sigma}\big)$.

The equality $(W_0(b) - \lambda)G(x-y) = \delta(x-y)$ could be interpreted as an equation of distributions.
Since the functions $f_{\pm}(x,k)$ are Jost solutions, the distribution defined by $(W_0(b) - \lambda)\times G(x-y)$ is supported only at $x = y$, and its singular part coincides with the singular part of the distribution
\begin{gather*}
-\frac{2\sigma\sigma'}{\pi \ii C(f_-,f_+)(\varkappa)} \bigg(\frac{f_-(x+2\sigma',\varkappa) f_+(y,\varkappa)- f_-(y,\varkappa)f_+(x+2\sigma',\varkappa)}{x-y-\ii0}
\\ \hphantom{-\frac{2\sigma\sigma'}{\pi \ii C(f_-,f_+)(\varkappa)} \bigg(}
{} + \frac{f_-(x-2\sigma',\varkappa) f_+(y,\varkappa)
- f_-(y,\varkappa)f_+(x-2\sigma',\varkappa)}{x-y+\ii0}\bigg)
\end{gather*}
in the neighbourhood of $x = y$. This singular part is equal to
\begin{gather*}
-\frac{2\sigma\sigma'}{\pi \ii} \bigg(\frac{1}{x-y-\ii0} - \frac{1}{x-y+\ii0} \bigg) = \delta(x-y),
\end{gather*}
where the authors used the Sokhotski--Plemelj formula. This formula is similar to the respective formula for a Schr\"odinger operator when the Dirac $\delta$-function appears by differentiating a step function.

\section[Proof of inequality (1.2)]{Proof of inequality (\ref{main})}

\subsection{Some auxiliary results}
We first collect some results from \cite{HLW} verbatim.
Let $A$ be a compact operator on a Hilbert space~$\mathcal{G}$ and let us denote
\begin{gather*}
\|A\|_n=\sum_{j=1}^n\sqrt{\lambda_j(A^*A)},
\end{gather*}
where $\lambda_j(A^*A)$ are the eigenvalues of $A^*A$ in decreasing order.
Then by Ky Fan's inequality (see for example \cite[Lemma~4.2]{GK}) the functionals $\|\cdot\|_n$, $n = 1,2,\dots $,
are norms and thus for any unitary operator $Y$ in $\mathcal{G}$ we have
\begin{gather*}
\|Y^* A Y\|_n = \|A\|_n.
\end{gather*}

\begin{Definition}
Let $A$, $B$ be two compact operators on $\mathcal{G}$. We say that $A$ majorises $B$ or $B\prec A$, iff
\begin{gather*}
\|B\|_n \le \|A\|_n, \qquad \text{for all} \quad n \in \mathbb N.
\end{gather*}
\end{Definition}

\begin{Lemma}\label{lem:maj}
Let $A$ be a nonnegative compact operator acting in $\mathcal{G}$, $\{Y(k)\}_{k\in \mathbb R}$ be a family of unitary operators on $\mathcal{G}$, and let $g(k)\, {\rm d} k$ be a probability measure on $\mathbb R$. Then the operator
\begin{gather*}
B = \int_{\mathbb R} Y(k)^* A Y(k) g(k) \, {\rm d} k
\end{gather*}
is majorised by~$A$.
\end{Lemma}

\begin{proof}This is a simple consequence of the triangle inequality
\begin{gather*}
\|B\|_n \le \int_{\mathbb R} \| Y^*(k) A Y(k)\|_n g(k) \, {\rm d} k = \|A\|_n \int_{\mathbb R} g(k) \, {\rm d} k = \|A\|_n.\tag*{\qed}
\end{gather*}
 \renewcommand{\qed}{}
\end{proof}

Let $\lambda_j=-2\cos\omega_j\le 2$ be the eigenvalues of $W_0(b)-V$ with $V\ge0$. In order to slightly simplify the notations it is convenient to write
\begin{gather*}
\lambda_j=-2\cos\big(\sqrt{\theta_j}\big)
\end{gather*}
with $\theta_j\in\big({-}\infty,\pi^2\big]$ and $\omega_j^2=\theta_j$.

Let us denote by $K_\lambda$ the Birman--Schwinger operator
\begin{gather}\label{BSch}
K_\lambda = V^{1/2} G_\lambda V^{1/2}.
\end{gather}
Let $\mu_j(K_\lambda)$ be the eigenvalues (in decreasing order) of the Birman--Schwinger operator $K_\lambda$ defined in \eqref{BSch}.
Then due to the Birman--Schwinger principle we have
\begin{gather}\label{BSch1}
1 = \mu_j(K_{\lambda_j}).
\end{gather}
Let us define the operator
\begin{gather*}%\label{Ltheta}
L_{\theta}:=\frac{1}{G_{-2\cos\sqrt{\theta}}(0)} K_{-2\cos\sqrt{\theta}} ,
\end{gather*}
where $G_{-2\cos\sqrt{\theta}}(0) =\frac{1}{2\pi b}\frac{\sqrt{\theta}}{\sin\sqrt{\theta}}$ is given in \eqref{G0}.
Then from \eqref{BSch1} we obtain
\begin{gather*}
\sum_{j\ge 1} \frac{1}{G_{\lambda_j}(0)}
=\sum_{j\ge 1} \frac{1}{G_{\lambda_j}(0)} \mu_j(K_{\lambda_j})=\sum_{j\ge 1}\mu_j(L_{\theta_j}).
\end{gather*}
The integral kernel of the operator $L_\theta$ is given by $\sqrt{V(x)}g_{\pi^2,\theta}(x-y)\sqrt{V(y)}$, where
\begin{gather*}
g_{\pi^2,\theta}(x):=
\frac{\pi}{\sqrt{\theta}}\frac{\sinh \big(\frac{\sqrt{\theta}}{b}x\big)}{\sinh \big(\frac{\pi}{b}x\big)}.
\end{gather*}
Consider a more general function
\begin{gather*}
g_{\varphi,\theta}(x):=
\frac{\sqrt\varphi}{\sqrt{\theta}}\frac{\sinh \big(\frac{\sqrt{\theta}}{b}x\big)}{\sinh \big(\frac{\sqrt\varphi}{b}x\big)}.
\end{gather*}
Since $g_{\varphi,\theta}(0)=1$ its Fourier transform
$\widehat{g}_{\varphi,\theta} = \mathcal{F} (g_{\varphi,\theta})$ satisfies the equation
\begin{gather*}
\int_{\mathbb R} \widehat{g}_{\varphi,\theta}(k) \, {\rm d} k = 1.
\end{gather*}
Moreover, for any $-\infty<\theta <\varphi$ with $0<\varphi <\pi^2$ we have
\begin{gather*}
 \widehat{g}_{\varphi,\theta}(k) = \mathcal{F} \Bigg(\frac{\sqrt\varphi}{\sqrt{\theta}} \frac{\sinh\big(\frac{\sqrt{\theta}}{b}x\big)}{\sinh\big(\frac{\sqrt\varphi}{b}x\big)}\Bigg)(k)
 = \frac{2\pi \sin\big(\pi \frac{\sqrt\theta}{\sqrt\varphi}\big)}{\sqrt\theta} \frac{b}{2\cosh \big(\frac{2\pi^2 b k}{\sqrt\varphi}\big) + 2\cos\big(\frac{\pi\sqrt\theta}{\sqrt\varphi}\big)},
\end{gather*}
and the right-hand side is positive. Thus $\widehat{g}_{\varphi,\theta}\, {\rm d} k$ is a probability measure for such values.

Note also that importantly
\begin{gather*}
\frac{g_{\pi^2,\theta}(x)}{g_{\pi^2,\theta'}(x)}
=\frac{\sqrt{\theta'}}{\sqrt{\theta}}\frac{\sinh\big(\frac{\sqrt{\theta}}{b}x\big)} {\sinh\big(\frac{\sqrt{\theta'}}{b}x\big)}=g_{\theta',\theta}(x)
\end{gather*}
and therefore
\begin{gather*}
\big(\widehat{g}_{\pi^2,\theta'}*\widehat{g}_{\theta',\theta}\big)(k)=\widehat{g}_{\pi^2,\theta}(k).
\end{gather*}
This is the interesting convolution/semigroup property mentioned in the introduction. In the special case $-\infty<\theta <0=\theta'$ analogous computations lead to the same result with $\widehat{g}_{0,\theta}(k)=\chi_{[-1,1]}\big(2\pi bk/\sqrt{|\theta|}\big)\pi b/\sqrt{|\theta|}$.

\begin{Lemma}[monotonicity]\label{Mon}
For $(\theta, \theta')$ such that $-\infty<\theta \le \theta'$ and $0\le\theta' <\pi^2$ we have $L_\theta\prec L_{\theta'}$.
\end{Lemma}

\begin{proof}
Let $Y(k)\colon L^2(\R)\to L^2(\R)$ be the unitary multiplication operator
\begin{gather*}
(Y(k)\psi)(x)=\e^{-2\pi \ii k x}\psi(x)
\end{gather*}
and let $T$ be the projection onto $V^{1/2}$, i.e.,
\begin{gather*}
(T\psi)(x)=V^{1/2} (x) \int_\R V^{1/2}(y) \psi(y) \, {\rm d} y.
\end{gather*}
Using $Y(k'+k'')=Y(k')Y(k'')$ and Lemma \ref{lem:maj} we obtain
\begin{align*}
L_{\theta}&=\int_\R Y(k)^*T Y(k) \widehat{g}_{\pi^2,\theta}(k)\, {\rm d} k
\\
&=\int_\R\int_{\R} Y(k)^*T Y(k) \widehat{g}_{\pi^2,\theta'}(k') \widehat{g}_{\theta',\theta}(k-k')\, {\rm d} k' {\rm d} k
\\
&=\int_\R Y(k'')^*\bigg(\int_{\R} Y(k')^*T Y(k') \widehat{g}_{\pi^2,\theta'}(k')\, {\rm d} k'\bigg)Y(k'') \widehat{g}_{\theta',\theta}(k'')\, {\rm d} k'' \prec L_{\theta'},
\end{align*}
where we have used that $\widehat{g}_{\theta',\theta} \, {\rm d} k$ is a probability measure.
\end{proof}

\begin{Remark}
With a slight abuse of notations, Lemma \ref{Mon} says that $L_{\lambda}\prec L_{\lambda'}$ for any $\lambda<2$ as long as $\lambda\le \lambda'$ and $-2\le\lambda' <2$.
\end{Remark}

\subsection[Proof of inequality (1.2)]{Proof of inequality (\ref{main})}

We now enumerate the eigenvalues of the operator $W_V(b)$ belonging to the interval $[-2,2)$ such that
$-2\le \lambda_1\le\lambda_2\le\lambda_3\le \cdots$ repeated with multiplicity.
By using the monotonicity established in Lemma~\ref{Mon} we have a sequence of inequalities
\begin{gather*}
\frac{1}{G_{\lambda_1}(0)} =2\pi b \frac{\sin\omega_1}{\omega_1} = \mu_1(L_{\theta_1}) \le \mu_1(L_{\theta_2}),
\\
\sum_{j=1}^2 \frac{1}{G_{\lambda_j}(0)} = 2\pi b \sum_{j=1}^2 \frac{\sin\omega_j}{\omega_j}
\le \sum_{j=1}^2 \mu_j(L_{\theta_2})
\le \sum_{j=1}^2 \mu_j(L_{\theta_3}),
\\
\sum_{j=1}^3 \frac{1}{G_{\lambda_j}(0)} = 2\pi b \sum_{j=1}^3 \frac{\sin\omega_j}{\omega_j}
\le \sum_{j=1}^3 \mu_j(L_{\theta_3})
\le \sum_{j=1}^3 \mu_j(L_{\theta_4}), \qquad {\rm etc}.
\end{gather*}
Note that we do not use any assumptions on the multiplicities of the eigenvalues, other than their finiteness. Furthermore, by Lemma \ref{Mon} the same results also hold true if a single eigenvalue is below $-2$.
Continuing the above process and noting that the trace of $L_{\theta}$ is $\int_\R V \, {\rm d} x$ for all $\theta$, we finally obtain
\begin{gather*}
\sum_{j\ge1}\frac{\sin\omega_j}{\omega_j}\le \frac{1}{2\pi b}\int_{\R}V(x)\, {\rm d} x.
\end{gather*}
The proof is complete.
